% math0 manuscript. Build with: latexmk -pdf paper.tex
\documentclass[11pt]{amsart}
\usepackage[margin=1.1in]{geometry}
\usepackage{amsmath,amssymb,amsthm,mathtools}
\usepackage{url}
\usepackage{hyperref}
\usepackage{xcolor}
\usepackage{listings}
\usepackage{newunicodechar}

\lstdefinelanguage{lean}{
  keywords={theorem,lemma,def,abbrev,example,fun,by,have,show,from,let,in,if,then,else,match,with,
    namespace,end,open,import,section,variable,instance,structure,where,calc,obtain,intro,exact,
    refine,rcases,rintro,simp,omega,decide,norm_num,linarith,nlinarith,positivity,induction,cases,
    private,noncomputable,rw,constructor,by_contra,push_cast,exact_mod_cast},
  sensitive=true, comment=[l]{--}, morecomment=[s]{/-}{-/}, string=[b]"}
\lstset{language=lean, basicstyle=\ttfamily\scriptsize, columns=fullflexible, keepspaces=true, breaklines=true,
  keywordstyle=\color{blue!60!black}\bfseries, commentstyle=\color{gray}, frame=single,
  extendedchars=true, inputencoding=utf8,
  literate={ℕ}{{$\mathbb{N}$}}1 {ℤ}{{$\mathbb{Z}$}}1 {∀}{{$\forall$}}1 {∃}{{$\exists$}}1 {∈}{{$\in$}}1
    {∧}{{$\wedge$}}1 {∨}{{$\vee$}}1 {¬}{{$\neg$}}1 {→}{{$\to$}}1 {←}{{$\leftarrow$}}1
    {↔}{{$\leftrightarrow$}}1 {≤}{{$\le$}}1 {≥}{{$\ge$}}1 {≠}{{$\ne$}}1 {≡}{{$\equiv$}}1 {∣}{{$\mid$}}1
    {×}{{$\times$}}1 {⟨}{{$\langle$}}1 {⟩}{{$\rangle$}}1 {·}{{$\cdot$}}1}
% Unicode commonly found in Lean sources.
\newunicodechar{ℕ}{\ensuremath{\mathbb{N}}}
\newunicodechar{ℤ}{\ensuremath{\mathbb{Z}}}
\newunicodechar{ℚ}{\ensuremath{\mathbb{Q}}}
\newunicodechar{ℝ}{\ensuremath{\mathbb{R}}}
\newunicodechar{ℂ}{\ensuremath{\mathbb{C}}}
\newunicodechar{∀}{\ensuremath{\forall}}
\newunicodechar{∃}{\ensuremath{\exists}}
\newunicodechar{→}{\ensuremath{\to}}
\newunicodechar{←}{\ensuremath{\leftarrow}}
\newunicodechar{↔}{\ensuremath{\leftrightarrow}}
\newunicodechar{↦}{\ensuremath{\mapsto}}
\newunicodechar{¬}{\ensuremath{\neg}}
\newunicodechar{∧}{\ensuremath{\wedge}}
\newunicodechar{∨}{\ensuremath{\vee}}
\newunicodechar{≤}{\ensuremath{\le}}
\newunicodechar{≥}{\ensuremath{\ge}}
\newunicodechar{≠}{\ensuremath{\ne}}
\newunicodechar{∈}{\ensuremath{\in}}
\newunicodechar{∉}{\ensuremath{\notin}}
\newunicodechar{⊆}{\ensuremath{\subseteq}}
\newunicodechar{∩}{\ensuremath{\cap}}
\newunicodechar{∪}{\ensuremath{\cup}}
\newunicodechar{∑}{\ensuremath{\sum}}
\newunicodechar{∏}{\ensuremath{\prod}}
\newunicodechar{⟨}{\ensuremath{\langle}}
\newunicodechar{⟩}{\ensuremath{\rangle}}
\newunicodechar{·}{\ensuremath{\cdot}}
\newunicodechar{∘}{\ensuremath{\circ}}
\newunicodechar{⁻}{\ensuremath{^{-}}}
\newunicodechar{¹}{\ensuremath{^{1}}}
\newunicodechar{²}{\ensuremath{^{2}}}
\newunicodechar{₀}{\ensuremath{_{0}}}
\newunicodechar{₁}{\ensuremath{_{1}}}
\newunicodechar{₂}{\ensuremath{_{2}}}
\newunicodechar{α}{\ensuremath{\alpha}}
\newunicodechar{β}{\ensuremath{\beta}}
\newunicodechar{ε}{\ensuremath{\varepsilon}}
\newunicodechar{δ}{\ensuremath{\delta}}
\newunicodechar{π}{\ensuremath{\pi}}
\newunicodechar{σ}{\ensuremath{\sigma}}
\newunicodechar{φ}{\ensuremath{\varphi}}
\newunicodechar{∞}{\ensuremath{\infty}}
\newunicodechar{⌊}{\ensuremath{\lfloor}}
\newunicodechar{⌋}{\ensuremath{\rfloor}}
\newunicodechar{⌈}{\ensuremath{\lceil}}
\newunicodechar{⌉}{\ensuremath{\rceil}}
\newunicodechar{∣}{\ensuremath{\mid}}
\newunicodechar{≡}{\ensuremath{\equiv}}
\newunicodechar{×}{\ensuremath{\times}}
\newunicodechar{⧸}{\ensuremath{/}}
\newunicodechar{‖}{\ensuremath{\|}}
\newunicodechar{√}{\ensuremath{\surd}}
\newunicodechar{∅}{\ensuremath{\emptyset}}
\newunicodechar{⦃}{\ensuremath{\{\!|}}
\newunicodechar{⦄}{\ensuremath{|\!\}}}
\newunicodechar{𝓝}{\ensuremath{\mathcal{N}}}
\newunicodechar{≫}{\ensuremath{\gg}}
\newunicodechar{≪}{\ensuremath{\ll}}

\theoremstyle{plain}
\newtheorem{theorem}{Theorem}
\newtheorem{lemma}[theorem]{Lemma}
\newtheorem{proposition}[theorem]{Proposition}
\newtheorem{corollary}[theorem]{Corollary}
\theoremstyle{definition}
\newtheorem{definition}[theorem]{Definition}
\newtheorem{remark}[theorem]{Remark}

\emergencystretch=3em
% Lean identifiers in running text (breakable at dots and underscores).
\expandafter\def\expandafter\UrlBreaks\expandafter{\UrlBreaks\do\_\do\.}
\DeclareUrlCommand\decl{\urlstyle{tt}}
\newcommand{\F}{\mathbb{F}_2}
\newcommand{\N}{\mathbb{N}}
\newcommand{\Z}{\mathbb{Z}}
\newcommand{\tsub}{\mathbin{\dot{-}}}
\newcommand{\co}[2]{[X^{#1}]\,#2}

\title[Sigrist's conjecture for OEIS A355339]{The central block of even entries in row $n$ of the knight-move triangle has length one exactly when $n=2^m-2$ with $m\ge2$: a proof of Sigrist's conjecture for OEIS A355339}
\author{math0}
\address{Produced by Claude (an AI model developed by Anthropic) in the math0 project}
\date{October 19, 2026}


\begin{document}
\begin{abstract}
Let $T(n,k)$ be the number of lattice paths from $(0,0)$ to $(n,k)$ with steps $(1,\pm2)$ and $(2,\pm1)$, so that row $n$ of the triangle A355320 of the OEIS is $T(n,-2n),\dots,T(n,2n)$, and let $a(n)$ (A355339) be the length of the maximal block of consecutive even entries of row $n$ containing $T(n,0)$, with $a(0)=0$.
We prove the conjecture recorded by R\'emy Sigrist in A355339 in 2022 that $a(n)=1$ if and only if $n=2^m-2$ for some $m\ge2$.
To do so we show that $T(n,0)$ is even for every $n\ge1$ and that $T(n,1)$ is odd exactly when $n+2$ is a power of $2$ at least $4$, using a closed form for the row polynomials over $\F$ and Lucas's theorem; to our knowledge this is the first proof of the conjecture.
The theorem is verified in Lean~4 against a statement that math0 formalized from the OEIS entry, reviewed for fidelity and frozen before the proof attempt, whose definitions are checked by the kernel against the $91$ listed terms of A355320 and the $77$ listed terms of A355339.
\end{abstract}
\maketitle

\section{Introduction}

A \emph{Catalan knight} in the sense of Kimberling \cite{A096587} moves from $(x,y)$ to one of $(x+1,y\pm2)$, $(x+2,y\pm1)$. For $n\ge0$ and $k\in\Z$ let $T(n,k)$ be the number of Catalan knight paths from $(0,0)$ to $(n,k)$; in the terminology of Baril, Hassler, Kirgizov and Ram\'irez \cite{BHKR} these are the grand knight's paths of size $n$ and altitude $k$. Since $|y|\le 2x$ along such a path, $T(n,k)=0$ for $|k|>2n$, and decomposing a path according to its last step gives $T(0,0)=1$ and
\begin{equation}\label{eq:rec}
  T(n,k)=T(n-1,k-2)+T(n-1,k+2)+T(n-2,k-1)+T(n-2,k+1)\qquad(n\ge1),
\end{equation}
where $T(-1,\cdot)=0$. Kimberling's triangle A096608 \cite{A096608} (2004) lists the right halves $0\le k\le 2n$ of the rows, and Sigrist's triangle A355320 \cite{A355320} (June~28, 2022) lists the full rows $-2n\le k\le 2n$, defined directly by \eqref{eq:rec}.

On June~29, 2022, Sigrist contributed A355339 \cite{A355339}: $a(n)$ is the number of central even terms in row $n$ of A355320, that is, the size of the maximal block of consecutive even entries $T(n,k)$ containing $k=0$ (and $0$ if $T(n,0)$ is odd), with $a(0)=0$. For example, row $4$ is $1,0,0,3,4,3,1,\mathbf{6},\mathbf{8},\mathbf{6},1,3,4,3,0,0,1$, so $a(4)=3$.
The entry conjectures that $a(n)=1$ if and only if $n>0$ and $n$ belongs to A000918 \cite{A000918}, the sequence $2^m-2$; that is, the central block is a single entry exactly for $n=2,6,14,30,62,\dots$. We prove this (Theorem~\ref{thm:main}). The proof rests on the following parity facts.

\begin{theorem}\label{thm:parity}
For all $n\ge0$ and $k\in\Z$: \textup{(1)} $T(n,-k)=T(n,k)$; \textup{(2)} if $n\ge1$, then $T(n,0)$ is even; \textup{(3)} $T(n,1)$ is odd if and only if $n+2=2^{j}$ for some $j\ge2$.
\end{theorem}

Part (1) is recorded in A355320, and (2), (3) concern $T(n,0)=\text{A096609}(n)$ and $T(n,1)=\text{A096610}(n)$. By (2) and (1), for $n\ge1$ the central block is a single entry if and only if $T(n,1)$ is odd.
To prove (2) and (3), we reduce row $n$ modulo $2$ to a polynomial $F_n\in\F[X]$ and show that $F_n=\sum_{p+q=n}\binom pq X^{3q}(1+X^2)^{2p-q}$. For $T(n,0)$ only central binomial coefficients occur, and these are even; for $T(n,1)$, Lucas's theorem \cite{Lucas} shows that only the term $p=q$ can be odd.

\subsection*{Status and prior work}
The statement is still marked ``Conjecture'' in A355339, whose current revision (\#10, June~29, 2022) dates from the day the entry was created (we read it in math0's copy of the OEIS database; the referees also checked the upstream OEIS data repository). The related entries A355320, A096587, A096608--A096612, A361292, A000918 and A002605 contain no parity results for $T(n,0)$ or $T(n,1)$.
Baril and Ram\'irez \cite{BR} showed that zigzag knight's paths (whose steps alternate up and down) of a given length ending on the $x$-axis are counted by Catalan numbers, and Baril, Hassler, Kirgizov and Ram\'irez \cite{BHKR} studied grand zigzag knight's paths. The only parity statement in \cite{BHKR} that we are aware of, Remark~3.2, says that the number of grand zigzag knight's paths ending at $(n,k)\ne(0,0)$ with $n\equiv k\pmod2$ is even; zigzag paths form a proper subfamily of the paths counted by $T(n,k)$, so this does not bear on the conjecture.
The referees' web searches (arXiv, MathOverflow, Mathematics Stack Exchange, the SeqFan list and general search engines) found no proof or counterexample, and the statement has no file in the formal-conjectures repository \cite{FC}.
We therefore classify the result as new. The argument is elementary, and we cannot exclude that it was known privately.

\section{Statement}\label{sec:statement}

\begin{definition}\label{def:T}
Define $T\colon\N\times\Z\to\N$ by $T(0,0)=1$, $T(n,k)=0$ for $|k|>2n$, and \eqref{eq:rec} for $n\ge1$, $|k|\le2n$, with $T(-1,\cdot)=0$. (Then \eqref{eq:rec} holds for all $n\ge1$, $k\in\Z$.)
For $n\ge1$ let $R(n)$ be the largest $r\in\{0,\dots,2n+1\}$ such that $T(n,0),T(n,1),\dots,T(n,r-1)$ are all even, and let $L(n)$ be defined in the same way with $T(n,0),T(n,-1),\dots,T(n,-(r-1))$. Put $a(0)=0$ and, for $n\ge1$, $a(n)=0$ if $R(n)=0$ and $a(n)=L(n)+R(n)-1$ otherwise.
\end{definition}

If $T(n,0)$ is even, $a(n)$ is the size of the maximal block $\{-(L(n)-1),\dots,R(n)-1\}$ of consecutive positions around $0$ with even entries.
The name of A355339 \cite{A355339} reads, verbatim, ``\texttt{a(n) is the number of central even terms in the n-th row of triangle A355320; a(0) = 0.}'', and its FORMULA section, in the entry as contributed by R\'emy Sigrist on June~29, 2022, reads, verbatim:
\begin{quote}
{\small\texttt{Conjecture: a(n) = 1 iff n > 0 and n belongs to A000918.}}
\end{quote}
The name of A355320 \cite{A355320} gives $T(0,0)=1$, and its FORMULA section states \eqref{eq:rec} for $n>0$, with $T$ extended by $0$ outside $-2n\le k\le 2n$; this is Definition~\ref{def:T}. The sequence A000918 \cite{A000918} is ``\texttt{a(n) = 2\^{}n - 2}'' with offset $0$, with terms $-1,0,2,6,14,\dots$.

\begin{theorem}[Sigrist's conjecture]\label{thm:main}
For every $n\ge0$, $a(n)=1$ if and only if $n>0$ and $n=2^m-2$ for some integer $m\ge0$.
\end{theorem}

\subsection*{Lean formalization}
The statement was formalized by math0 in \path{lean/Math0/Statements/OEIS/A355339.lean}, reviewed for fidelity and frozen (SHA-256 \texttt{ea5b1be9\ldots 4226b}) before any proof attempt. Its definitions and the statement read, verbatim (the import, comments and docstrings, blank lines and all theorems are omitted):
\begin{lstlisting}
namespace Math0.Statements.A355339
def addL : List ℕ → List ℕ → List ℕ
  | [], l => l
  | x :: xs, [] => x :: xs
  | x :: xs, y :: ys => (x + y) :: addL xs ys
def step (pq : List ℕ × List ℕ) : List ℕ × List ℕ :=
  (pq.2, addL (addL pq.2 (List.replicate 4 0 ++ pq.2))
    (addL (List.replicate 3 0 ++ pq.1) (List.replicate 5 0 ++ pq.1)))
def rowPair (n : ℕ) : List ℕ × List ℕ := step^[n] ([], [1])
def row (n : ℕ) : List ℕ := (rowPair n).2
def entry (r : List ℕ) (m : ℕ) (k : ℤ) : ℕ :=
  if 0 ≤ k + 2 * m then r.getD (k + 2 * m).toNat 0 else 0
def T (n : ℕ) (k : ℤ) : ℕ := entry (row n) n k
def a (n : ℕ) : ℕ :=
  if n = 0 then 0 else
    let r := row n
    let right := (((List.range (2 * n + 1)).map fun k : ℕ => entry r n k).takeWhile
      fun x => x % 2 == 0).length
    let left := (((List.range (2 * n + 1)).map fun k : ℕ => entry r n (-k)).takeWhile
      fun x => x % 2 == 0).length
    if right = 0 then 0 else left + right - 1
def statement : Prop :=
  ∀ n : ℕ, a n = 1 ↔ (0 < n ∧ ∃ m : ℕ, (n : ℤ) = 2 ^ m - 2)
end Math0.Statements.A355339
\end{lstlisting}
Row $n$ is stored as the list \texttt{row n} of its $4n+1$ entries, with $T(n,k)$ at position $k+2n$; \texttt{entry r n k} reads position $k+2n$ of \texttt{r} and returns $0$ outside the list. The map \texttt{step} sends (row $n-1$, row $n$) to (row $n$, row $n+1$), starting from row $-1$ (empty) and row $0$. Position $k+2n+2$ of row $n+1$ receives positions $k+2n+2$ and $k+2n-2$ of row $n$ (list shifts $0$ and $4$), that is $T(n,k+2)$ and $T(n,k-2)$, and positions $k+2n-1$ and $k+2n-3$ of row $n-1$ (shifts $3$ and $5$), that is $T(n-1,k+1)$ and $T(n-1,k-1)$; this is the recurrence $R_{n+1}=(1+X^4)R_n+(X^3+X^5)R_{n-1}$ of the PARI program in A355320. In \decl{a}, \texttt{right} and \texttt{left} are $R(n)$ and $L(n)$ (\texttt{takeWhile} takes the longest prefix of even values). Membership in A000918 is written in $\Z$, so $2^m-2$ is not truncated, and the statement covers all $n\ge0$ without further hypotheses.

The statement file proves that \decl{T} satisfies the defining equations of Definition~\ref{def:T}, which determine it by induction on $n$: \decl{T_zero_zero} ($T(0,0)=1$), \decl{T_eq_zero} ($T(n,k)=0$ for $2n<|k|$), \decl{T_one} ($T(1,k)=T(0,k-2)+T(0,k+2)$) and \decl{T_add_two} (\eqref{eq:rec} for $T(n+2,k)$, every $k\in\Z$). The theorem \decl{a_eq} restates \decl{a} with \texttt{T n k} in place of \texttt{entry r n k} (by \texttt{rfl}). The theorems \decl{T_data} and \decl{a_data}, proved by \texttt{decide +kernel}, state that the values $T(n,k)$ for $0\le n\le 6$ are the $91$ listed terms of A355320 and that $a(0),\dots,a(76)$ are the $77$ listed terms of A355339. Independent Python computations by the reviewers agree with Sigrist's b-file of A355339 for $n\le 8192$ (they are not part of the formal proof).

\section{Proof}\label{sec:proof}

We follow the Lean file \path{lean/Math0/Frontier/OEIS/A355339.lean}; each step names the corresponding declaration. Throughout, $\binom pq=0$ for $q>p$, $x\tsub y=\max(x-y,0)$ on $\N$, and $\co{m}{P}$ is the coefficient of $X^m$ in $P$. All polynomials are in $\F[X]$, and an integer is identified with its image in $\F$.

\subsection{\texorpdfstring{Row $0$ and symmetry}{Row 0 and symmetry}}

\begin{lemma}[\decl{T0}, \decl{T_symm}]\label{lem:symm}
$T(0,k)=[k=0]$, and $T(n,-k)=T(n,k)$ for all $n\ge0$, $k\in\Z$.
\end{lemma}

\begin{proof}
The first claim is \decl{T_zero_zero} and \decl{T_eq_zero}. Symmetry follows by two-step induction on $n$: rows $0$ and $1$ are symmetric ($T(1,k)=[k=\pm2]$), and by the induction hypothesis the four terms $T(n+1,-k\mp2)$, $T(n,-k\mp1)$ of \eqref{eq:rec} for $T(n+2,-k)$ equal the four terms $T(n+1,k\pm2)$, $T(n,k\pm1)$ for $T(n+2,k)$.
\end{proof}

\subsection{\texorpdfstring{Row polynomials over $\mathbb{F}_2$}{Row polynomials over F2}}
Define $F_n\in\F[X]$ (\decl{Fp}) by $F_0=1$, $F_1=1+X^4$ and $F_{n+2}=(1+X^4)F_{n+1}+(X^3+X^5)F_n$.

\begin{lemma}[\decl{T_shift}, \decl{T_coeff}]\label{lem:coeff}
For all $n,m\in\N$, $T(n,m-2n)\equiv\co{m}{F_n}\pmod 2$.
\end{lemma}

\begin{proof}
For $j,m\in\N$ we have $T(n,m-2n-j)=T\bigl(n,(m-j)-2n\bigr)$ if $j\le m$, and $T(n,m-2n-j)=0$ if $j>m$, since then $m-2n-j<-2n$ (\decl{T_shift}).
We use two-step induction on $n$. For $n=0$, $T(0,m)=[m=0]=\co{m}{1}$. For $n=1$, $T(1,m-2)=T(0,m-4)+T(0,m)=[m=4]+[m=0]=\co{m}{(1+X^4)}$. For $n+2$, put $k=m-2(n+2)$; then $k-2=m-2(n+1)-4$, $k+2=m-2(n+1)$, $k-1=m-2n-5$ and $k+1=m-2n-3$, so by \eqref{eq:rec}, the shift rule and the induction hypothesis, modulo $2$,
\[
  T(n+2,k)\equiv \co{m}{F_{n+1}}+[m\ge4]\,\co{m-4}{F_{n+1}}+[m\ge5]\,\co{m-5}{F_n}+[m\ge3]\,\co{m-3}{F_n},
\]
which is the coefficient of $X^m$ in $(1+X^4)F_{n+1}+(X^3+X^5)F_n=F_{n+2}$.
\end{proof}

\subsection{\texorpdfstring{A closed form for $F_n$}{A closed form for F\_n}}

\begin{lemma}[\decl{hsum_rec}]\label{lem:hsum}
Let $S$ be a commutative semiring, $A,B\in S$, and $h_n=\sum_{p+q=n}\binom pq A^pB^q$ (\decl{hsum}). Then $h_0=1$, $h_1=A$ and $h_{n+2}=A\,h_{n+1}+AB\,h_n$ for all $n\ge0$.
\end{lemma}

\begin{proof}
The values of $h_0,h_1$ are immediate. In $h_{n+2}$ the term $p=0$ vanishes, so $h_{n+2}=\sum_{p+q=n+1}\binom{p+1}{q}A^{p+1}B^q$. By Pascal's rule, $\binom{p+1}{0}=\binom p0$ and $\binom{p+1}{q+1}=\binom pq+\binom p{q+1}$; hence
\[
  h_{n+2}=\sum_{p+q=n+1}\binom pqA^{p+1}B^q+\sum_{p+q=n}\binom pqA^{p+1}B^{q+1}=A\,h_{n+1}+AB\,h_n.\qedhere
\]
\end{proof}

\begin{lemma}[\decl{one_add_X4}, \decl{Fp_mul}, \decl{Fp_closed}]\label{lem:closed}
Let $v=1+X^2$. Then $F_nv^n=h_n(v^3,X^3)$ for all $n\ge0$, and
$F_n=\sum_{p+q=n}\binom pq X^{3q}\,v^{\,2p\tsub q}$.
\end{lemma}

\begin{proof}
In $\F[X]$ we have $1+X^4=(1+X^2)^2-2X^2=v^2$, and $X^3+X^5=X^3v$. So $G_n=F_nv^n$ satisfies $G_0=1$, $G_1=v^2\cdot v=v^3$ and
\[
  G_{n+2}=v^2F_{n+1}v^{n+2}+X^3vF_nv^{n+2}=v^3G_{n+1}+v^3X^3G_n,
\]
which are the initial values and the recurrence of $h_n(v^3,X^3)$ in Lemma~\ref{lem:hsum}.
For the closed form, a term with $q>p$ vanishes, and for $q\le p$ we have $3p=(2p-q)+n$, so $\binom pqX^{3q}v^{2p-q}\cdot v^n=\binom pq(v^3)^p(X^3)^q$. Hence both sides of the claimed identity, multiplied by $v^n$, equal $h_n(v^3,X^3)$, and $v^n$ can be cancelled because $\F[X]$ is an integral domain and $v\ne0$.
\end{proof}

\begin{lemma}[\decl{vpow_coeff}, \decl{Fp_coeff}]\label{lem:fcoeff}
For $e,j\in\N$, $\co{j}{v^e}=\binom{e}{j/2}$ if $j$ is even and $0$ if $j$ is odd. Hence, for all $n,m\in\N$, in $\F$,
\begin{equation}\label{eq:coeff}
  \co{m}{F_n}=\sum_{p+q=n}\binom pq\,\varepsilon_m(p,q),\qquad
  \varepsilon_m(p,q)=\begin{cases}\binom{2p\tsub q}{(m-3q)/2}&\text{if }3q\le m\text{ and }2\mid m-3q,\\ 0&\text{otherwise.}\end{cases}
\end{equation}
\end{lemma}

\begin{proof}
$v^e=\sum_i\binom eiX^{2i}$ (in Lean, $v^e$ is the image of $(1+X)^e$ under $X\mapsto X^2$, \decl{Polynomial.expand}). The formula follows from Lemma~\ref{lem:closed}, since $\co{m}{(X^{3q}P)}=\co{m-3q}{P}$ if $3q\le m$ and $0$ otherwise.
\end{proof}

\subsection{Two binomial parity lemmas}
We use Lucas's theorem \cite{Lucas} for the prime $2$ (\decl{lucas2}, from Mathlib's \decl{Choose.choose_modEq_choose_mod_mul_choose_div_nat}),
\begin{equation}\label{eq:lucas}
  \tbinom NK\equiv\tbinom{N\bmod2}{K\bmod2}\tbinom{\lfloor N/2\rfloor}{\lfloor K/2\rfloor}\pmod 2,
\end{equation}
and the fact that $\binom{2c}{c}=2\binom{2c-1}{c-1}$ is even for $c\ge1$ (\decl{Nat.two_dvd_centralBinom_of_one_le}).

\begin{lemma}[\decl{oddC}]\label{lem:oddC}
For $b\ge0$, $\binom{2b+1}{b}$ is odd if and only if $b+1$ is a power of $2$.
\end{lemma}

\begin{proof}
Strong induction on $b$. By \eqref{eq:lucas}, $\binom{2b+1}{b}\equiv\binom{1}{b\bmod2}\binom{b}{\lfloor b/2\rfloor}=\binom{b}{\lfloor b/2\rfloor}$.
If $b=0$, this is $1$ and $b+1=2^0$.
If $b=2c$ with $c\ge1$, this is $\binom{2c}{c}$, which is even, and $b+1\ge3$ is odd, so not a power of $2$.
If $b=2c+1$, this is $\binom{2c+1}{c}$, which by induction is odd if and only if $c+1=2^t$ for some $t$, that is, if and only if $b+1=2(c+1)$ is a power of $2$ (note $b+1\ne 2^0$).
\end{proof}

\begin{lemma}[\decl{dzero}]\label{lem:dzero}
Let $t,D,q\ge0$. If $\binom{D+q}{q}$ is odd and $2D+q+1=2^t$, then $D=0$.
\end{lemma}

\begin{proof}
Induction on $t$. If $t=0$, then $2D+q+1=1$ forces $D=0$. Let $2D+q+1=2^{t+1}$; then $q=2q'+1$ is odd. If $D$ were odd, $D+q$ would be even and \eqref{eq:lucas} would give the factor $\binom01=0$. Hence $D=2D'$, and \eqref{eq:lucas} gives $\binom{D+q}{q}\equiv\binom11\binom{D'+q'}{q'}$, so $\binom{D'+q'}{q'}$ is odd. Moreover $2(2D'+q'+1)=2D+q+1=2^{t+1}$, so $2D'+q'+1=2^t$, and the induction hypothesis gives $D'=0$.
\end{proof}

\subsection{\texorpdfstring{The entries $T(n,0)$ and $T(n,1)$}{The entries T(n,0) and T(n,1)}}

\begin{proposition}[\decl{term_even}, \decl{T0_even}]\label{prop:T0}
For $n\ge1$, $T(n,0)$ is even.
\end{proposition}

\begin{proof}
By Lemma~\ref{lem:coeff}, $T(n,0)\equiv\co{2n}{F_n}$; we show that every term of \eqref{eq:coeff} with $m=2n$ vanishes. Let $p+q=n$. If $p<q$, then $\binom pq=0$. Let $p\ge q$, and suppose $3q\le 2n$ and $2n-3q$ is even (otherwise $\varepsilon_{2n}(p,q)=0$). Since $2n-3q=2p-q$, we can write $2p-q=2c$, and $\varepsilon_{2n}(p,q)=\binom{2c}{c}$. As $n\ge1$ and $p\ge q$, we have $p\ge1$ and $2c=2p-q\ge p\ge1$, so $c\ge1$ and $\binom{2c}{c}$ is even.
\end{proof}

\begin{proposition}[\decl{term_off}, \decl{T1_parity}]\label{prop:T1}
For $n\ge0$, $T(n,1)$ is odd if and only if $n+2=2^{t+2}$ for some $t\ge0$.
\end{proposition}

\begin{proof}
By Lemma~\ref{lem:coeff}, $T(n,1)\equiv\co{2n+1}{F_n}$.

\emph{Off-diagonal terms vanish.} Let $p+q=n$ with $p\ne q$; we show $\binom pq\varepsilon_{2n+1}(p,q)=0$ in $\F$. If $p<q$, then $\binom pq=0$. Let $p>q$. Then $3q\le2n+1$, and $2n+1-3q=(2p-q)+1$. If this is odd, $\varepsilon_{2n+1}(p,q)=0$. Otherwise $2p-q=2c+1$ is odd and $\varepsilon_{2n+1}(p,q)=\binom{2c+1}{c+1}=\binom{2c+1}{c}$.
Suppose both $\binom pq$ and $\binom{2c+1}{c}$ were odd. By Lemma~\ref{lem:oddC}, $c+1=2^t$ for some $t$. Put $D=p-q\ge1$; then $\binom{D+q}{q}=\binom pq$ is odd and $2D+q+1=2p-q+1=2c+2=2^{t+1}$, so Lemma~\ref{lem:dzero} gives $D=0$, a contradiction.

\emph{$n$ odd.} Then $p+q=n$ forces $p\ne q$, so all terms vanish and $T(n,1)$ is even; and $n+2$ is odd, hence not of the form $2^{t+2}$.

\emph{$n=2q$ even.} Only the term $(p,q)=(q,q)$ remains; here $\binom qq=1$, $3q\le 4q+1$ and $2n+1-3q=q+1$.
If $q=2c$ is even, then $q+1$ is odd, the term is $0$ and $T(n,1)$ is even; and $n+2=4c+2\equiv2\pmod4$ is not of the form $2^{t+2}$.
If $q=2c+1$ is odd, then $(q+1)/2=c+1$ and $2q\tsub q=2c+1$, so $T(n,1)\equiv\binom{2c+1}{c+1}=\binom{2c+1}{c}$, which by Lemma~\ref{lem:oddC} is odd if and only if $c+1=2^t$ for some $t$, that is, if and only if $n+2=4(c+1)=2^{t+2}$.
\end{proof}

Lemma~\ref{lem:symm} and Propositions~\ref{prop:T0} and~\ref{prop:T1} prove Theorem~\ref{thm:parity}.

\subsection{\texorpdfstring{The run length and the proof of Theorem~\ref{thm:main}}{The run length and the proof of Theorem 3}}

\begin{lemma}[\decl{range_map}, \decl{tw}, \decl{run_one}, \decl{a_one_iff}]\label{lem:run}
For $n\ge1$, $a(n)=1$ if and only if $T(n,1)$ is odd.
\end{lemma}

\begin{proof}
Since $n\ge1$, the lists $T(n,0),T(n,1),\dots,T(n,2n)$ and $T(n,0),T(n,-1),\dots,T(n,-2n)$ have at least two entries, and their first entry $T(n,0)$ is even by Proposition~\ref{prop:T0}. So $R(n),L(n)\ge1$, $a(n)=L(n)+R(n)-1$, and $a(n)=1$ if and only if $L(n)=R(n)=1$. Finally, $R(n)=1$ if and only if $T(n,1)$ is odd, and $L(n)=1$ if and only if $T(n,-1)$ is odd, which is the same condition by Lemma~\ref{lem:symm}.
\end{proof}

\begin{proof}[Proof of Theorem~\ref{thm:main}]
For $n=0$ both sides are false, as $a(0)=0$. Let $n\ge1$. By Lemma~\ref{lem:run} and Proposition~\ref{prop:T1}, $a(n)=1$ if and only if $n+2=2^{t+2}$ for some $t\ge0$. If $n+2=2^{t+2}$, take $m=t+2$. Conversely, if $n=2^m-2$ in $\Z$ with $n>0$, then $m\ne0$ (else $n=-1$) and $m\ne1$ (else $n=0$), so $m=t+2$ and $n+2=2^{t+2}$.
\end{proof}

\section{Machine verification}\label{sec:lean}

\subsection*{Files and declarations}
The proof is in the file \path{lean/Math0/Frontier/OEIS/A355339.lean} of the math0 repository, \url{https://github.com/Foundation-0/math0}. It imports the frozen statement file \path{lean/Math0/Statements/OEIS/A355339.lean} (Section~\ref{sec:statement}) and the audit commands in \path{lean/Math0/Audit.lean}. The main theorem is \decl{Math0.Frontier.A355339.result}, of type \decl{Math0.Statements.A355339.statement}; the auxiliary declarations of the proof file named in Section~\ref{sec:proof} are private. The proof file uses no \texttt{decide}, \texttt{native\_decide}, \texttt{sorry} or \texttt{set\_option}. The data checks \decl{T_zero_zero}, \decl{T_data} and \decl{a_data} in the statement file use \texttt{decide +kernel}, which evaluates in the kernel and adds no axioms.

\subsection*{Audit}
The proof file ends with \texttt{\#math0\_settles result statement} (with full names). This command, defined in \path{lean/Math0/Audit.lean} for statements that math0 formalizes itself, fails unless (1) \decl{statement} is a definition of type \texttt{Prop} without universe parameters; (2) \decl{statement} and \decl{result} use no axioms besides \decl{propext}, \decl{Classical.choice} and \decl{Quot.sound} (so no \decl{sorryAx}, \texttt{native\_decide} or custom axiom); and (3) the type of \decl{result} is syntactically the statement or its negation, up to metadata annotations and without unfolding. Here it reports
\begin{lstlisting}[basicstyle=\ttfamily\tiny]
math0 audit: 'Math0.Statements.A355339.statement' depends on axioms [propext]
math0 audit: 'Math0.Frontier.A355339.result' depends on axioms [Classical.choice, Quot.sound, propext]
math0 statement check: 'Math0.Frontier.A355339.result' proves 'Math0.Statements.A355339.statement'
\end{lstlisting}

\subsection*{Toolchain and reproduction}
We use Lean~4 \cite{Lean4} \texttt{v4.33.1} and Mathlib \cite{Mathlib} \texttt{v4.33.1}, pinned in \path{lean/}; the repository also pins formal-conjectures \cite{FC} at commit \texttt{c77fdd4}, which this result does not use.
To reproduce, run \path{scripts/verify.sh} from the repository root: it fetches the Mathlib build cache, builds all math0 results with \texttt{lake build} (re-running every audit command), checks that the pinned \texttt{formal\_conjectures} and \texttt{mathlib} checkouts are unmodified, and reports the number of passed statement checks. Checking the proof file alone with \texttt{lake env lean} takes about $15$ seconds on our machine.

{\footnotesize\raggedright
\begin{thebibliography}{99}

\bibitem{BR}
J.-L.~Baril and J.~L.~Ram\'irez, \emph{Knight's paths towards Catalan numbers}, arXiv:2206.12087 [math.CO], 2022.

\bibitem{BHKR}
J.-L.~Baril, N.~Hassler, S.~Kirgizov, and J.~L.~Ram\'irez, \emph{Grand zigzag knight's paths}, Enumer. Comb. Appl. \textbf{5}:1 (2025), Article \#S2R6, \url{https://doi.org/10.54550/ECA2025V5S1R6}; arXiv:2402.04851 [math.CO].

\bibitem{FC}
M.~Firsching, P.~Lezeau, S.~Mercuri, M.~Z. Horv\'ath, Y.~Dillies, C.~S\"onne, E.~Wieser, F.~Zhang, T.~Hubert, B.~Ag\"uera y Arcas, and P.~Kohli, \emph{Formal Conjectures: An Open and Evolving Benchmark for Verified Discovery in Mathematics}, arXiv:2605.13171 [cs.AI], 2026.

\bibitem{Lucas}
\'E.~Lucas, \emph{Th\'eorie des fonctions num\'eriques simplement p\'eriodiques}, Amer. J. Math. \textbf{1} (1878), 184--240, 289--321.

\bibitem{Mathlib}
The mathlib Community, \emph{The Lean mathematical library}, in: Proc.\ 9th ACM SIGPLAN Int.\ Conf.\ on Certified Programs and Proofs (CPP 2020), 367--381.

\bibitem{Lean4}
L.~de Moura and S.~Ullrich, \emph{The Lean 4 theorem prover and programming language}, in: CADE 28, Lecture Notes in Comput. Sci. \textbf{12699}, Springer, 2021, 625--635.

\bibitem{A000918}
The OEIS Foundation Inc., \url{https://oeis.org/A000918}, entry by N.~J.~A. Sloane.

\bibitem{A096587}
The OEIS Foundation Inc., \url{https://oeis.org/A096587}, entry by C.~Kimberling (June~28, 2004).

\bibitem{A096608}
The OEIS Foundation Inc., \url{https://oeis.org/A096608}, entry by C.~Kimberling (June~29, 2004).

\bibitem{A355320}
The OEIS Foundation Inc., \url{https://oeis.org/A355320}, entry by R.~Sigrist (June~28, 2022).

\bibitem{A355339}
The OEIS Foundation Inc., \url{https://oeis.org/A355339}, entry and conjecture by R.~Sigrist (June~29, 2022), rev.~\#10.

\end{thebibliography}
}

\end{document}
